% §4 Convolution representation --- blueprint chapter 4 (blueprint/src/parts/04-representation.tex),
% statements and proofs of record verbatim. Drafted 2026-09-11.
%
% Contents and departures from the blueprint chapter:
%   - both statement nodes transcribed (lem:convolution-representation, lem:nonvanishing) under
%     their markers; both [T] and \leanok, so no page anchor is printed; the proof of
%     lem:nonvanishing keeps the blueprint's two forms of the nonvanishing step, the chain form
%     being the machine-checked one and the printed text saying so;
%   - the lead paragraph and "the device" are the blueprint's prose; the sentence pointing at
%     the uniform blur's transform is expanded with the formula (the draft had it);
%   - the two remarks are copied (rem:modulation-transfer, rem:wendel); rem:wendel is this
%     section's one relation-to-the-causal-theory remark, and it is reworded to say so in its
%     title --- the blueprint's "against the causal version" is the same comparison;
%   - the blueprint's chapter references become section references.
% Forward references: prop:kernel-regularity (§7), thm:increments-levy (§5), sec:cascade,
% sec:covariance, prop:no-positivity-no-classification (§3).

\section{Convolution representation}\label{sec:representation}

This section derives the form of the measurement operators from the axioms that do not
involve the cascade or covariance. The aim is Lemma~\ref{lem:convolution-representation}:
under (A1)--(A5) every operator $\Phis{s}{t}$ is a convolution by a symmetric probability
measure $\mu_{s,t}$ on $\RR$, so the apertures of \S\ref{sec:axioms} are measures. One
operator $\Phis{s}{t}$, the measurement that carries the record at scale $s$ to the record
at scale $t$, is called a \emph{stage} below, and $\mu_{s,t}$ is its kernel. The lemma has
a converse, that every symmetric probability measure defines an operator satisfying
(A1)--(A5), and the constructive half of the characterization in
\S\ref{sec:characterization} builds its families through that converse. Whether the kernels
have densities is settled later, once covariance and nondegeneracy are imposed
(Proposition~\ref{prop:kernel-regularity}).

Read at the kernel, the lemma gives the picture the rest of the paper uses. A stage
reports the signal at a randomly displaced position, its kernel is the law of that
displacement, and the transform of the kernel is the stage's modulation transfer function
(Remark~\ref{rem:modulation-transfer}). The second result of the section,
Lemma~\ref{lem:nonvanishing}, adds the cascade (A6) and continuity (A7), drops positivity,
and shows that the transforms of the kernels do not vanish. That is what lets the rest of
the article work with one function of each kernel, its exponent $-\log\hat\mu_{s,t}$.

\emph{The device.} The characters $e_\omega(x) := e^{-i\omega x}$ are bounded on all of $\RR$,
so $e_\omega \in L^\infty$ and complex-valued functionals on the real space $\Lone$ may be used
freely, $\langle w, f\rangle := \int_\RR w f$. Equivalently one pairs with $\cos\omega x$ and
$\sin\omega x$. For a finite measure $\mu$ and $f \in \Lone$, Fubini gives
\begin{equation}\label{eq:pairing}
  \langle e_\omega, \mu * f\rangle = \hat\mu(\omega)\,\hat f(\omega), \qquad \omega \in \RR .
\end{equation}
With $f$ the standard Gaussian density, $\hat f(\omega) = e^{-\omega^2/2} > 0$, so a single
test function reads off the transform of the kernel at every frequency.

% shared with blueprint lem:convolution-representation
\begin{lemma}[representation]\label{lem:convolution-representation}
Under (A1), (A2), (A4), (A5), for each pair $0 \le s \le t$ there is a unique probability
measure $\mu_{s,t}$ on $\RR$ with
\[
  \Phis{s}{t} f = \mu_{s,t} * f, \qquad f \in \Lone ;
\]
under (A3) it is symmetric, $R\mu_{s,t} = \mu_{s,t}$. Conversely every symmetric probability
measure on $\RR$ defines an operator satisfying (A1)--(A5).
\end{lemma}

The lemma states symmetry as a separate clause, under (A3), rather than assuming
(A1)--(A5) at once, because the representation is also of interest without positivity.
Proposition~\ref{prop:no-positivity-no-classification} needs a representation of
operators satisfying (A1)--(A3) and (A5) but not (A4). That is not what the lemma gives, in
either of its forms; it is what Wendel's theorem gives (Remark~\ref{rem:wendel}).

\begin{proof}
Fix $(s,t)$ and write $\Phi = \Phis{s}{t}$.

\emph{Convolution as a vector-valued integral.} For $f, g \in \Lone$,
$f * g = \int_\RR f(y)\,\Tr{y}g\,dy$ as a Bochner integral in $\Lone$: the integrand is
strongly measurable, since $y \mapsto \Tr{y}g$ is continuous because translation acts
continuously on $\Lone$, and it is dominated by $|f(y)|\,\|g\|_1$. An $\Lone$-valued integral commutes with every bounded operator, so (A1) and
(A2) give the interchange identity $\Phi(f * g) = f * \Phi g$.

\emph{The approximants.} Let $\rho_\varepsilon := (2\varepsilon)^{-1}\mathbf 1_{(-\varepsilon,\varepsilon)}$
and $h_\varepsilon := \Phi\rho_\varepsilon$. By (A4) and (A5), $h_\varepsilon \ge 0$ with
$\int h_\varepsilon = 1$, so $\nu_\varepsilon := h_\varepsilon\,dx$ is a probability measure on
$\RR$.

\emph{Tightness.} By the interchange identity $\rho_1 * h_\varepsilon = \Phi(\rho_1 *
\rho_\varepsilon)$, and $\rho_1 * \rho_\varepsilon \to \rho_1$ in $\Lone$ as
$\varepsilon \downarrow 0$, so $\rho_1 * h_\varepsilon$ converges in $\Lone$; and a convergent
sequence in $\Lone$ has uniformly small tails. Convolution with a nonnegative unit-mass density
carried by $[-1,1]$ moves mass by at most $1$: for $|y| > R+1$ the translate
$\rho_1(\cdot - y)$ is carried by $\{|x| > R\}$, so by Tonelli
\[
  \int_{|x| > R} (\rho_1 * h_\varepsilon)(x)\,dx
  \;=\; \int_\RR h_\varepsilon(y)\int_{|x|>R}\rho_1(x-y)\,dx\,dy
  \;\ge\; \int_{|y| > R+1} h_\varepsilon(y)\,dy
  \;=\; \nu_\varepsilon(\{|y| > R+1\}).
\]
Along $\varepsilon_n = 1/(n+1)$ the tails of $\nu_{\varepsilon_n}$ are therefore small
uniformly in $n$: the family is tight, with compacts $[-R-1, R+1]$.

\emph{The limit.} By Prokhorov's theorem a subsequence $\nu_{\varepsilon_{n_k}}$ converges
weakly to a probability measure $\mu$ on $\RR$.

\emph{Identification.} Let $\Psi \in L^\infty$ and $f \in \Lone$. The map
$y \mapsto \Psi(\Tr{y}f)$ is continuous and bounded by $\|\Psi\|\,\|f\|_1$, hence a legitimate
test function for weak convergence. Reading $f * h_\varepsilon = \int h_\varepsilon(y)\,\Tr{y}f\,dy$
with $h_\varepsilon$ in the density slot,
\[
  \Psi(f * h_{\varepsilon_{n_k}}) = \int_\RR \Psi(\Tr{y}f)\,\nu_{\varepsilon_{n_k}}(dy)
  \ \xrightarrow{\ k \to \infty\ }\ \int_\RR \Psi(\Tr{y}f)\,\mu(dy) = \Psi(\mu * f),
\]
the last equality because $\mu * f = \int \Tr{y}f\,\mu(dy)$ is again an $\Lone$-valued Bochner
integral. On the other hand $f * h_\varepsilon = \Phi(f * \rho_\varepsilon) \to \Phi f$ in
$\Lone$, by the interchange identity and (A1). Hence $\Psi(\Phi f) = \Psi(\mu * f)$ for every
$\Psi \in L^\infty$, and $\Phi f = \mu * f$ for every $f \in \Lone$ at once.

\emph{Uniqueness.} Take $f$ the standard Gaussian density and pair against $e_\omega$: by
\ref{eq:pairing}, $\langle e_\omega, \Phi f\rangle = \hat\mu(\omega)\,e^{-\omega^2/2}$, so
$\Phi$ determines $\hat\mu$ on all of $\RR$, and a finite measure on $\RR$ is determined by its
transform (Proposition~\ref{prop:fourier-uniqueness}). One test function suffices.

\emph{Symmetry.} Under (A3), $R(\mu * f) = (R\mu) * (Rf)$ gives
$(R\mu) * f = R\,\Phi\,R f = \Phi f = \mu * f$ for every $f$, so $R\mu = \mu$ by uniqueness.

\emph{Converse.} A probability measure $\mu$ on $\RR$ makes $f \mapsto \mu * f$ bounded with
norm at most $1$ by Young's inequality (A1), translation-commuting because convolution is
(A2), positive (A4), and mass-preserving by Tonelli (A5); if $\mu$ is symmetric it commutes
with $R$ (A3).
\end{proof}

The rest of the article works with one function of the measures, the \emph{exponent}
$g_{s,t}(\omega) := -\log\hat\mu_{s,t}(\omega)$, the quantity the cascade makes additive
(\S\ref{sec:cascade}). For it to exist the transform must be positive, which on the line is
not automatic: a symmetric probability measure can have a negative or vanishing transform,
as the uniform density on $[-1,1]$ does, with $\hat\mu(\omega) = \sin\omega/\omega$.
Excluding this stage by stage needs the cascade and continuity beyond (A1)--(A5).

% shared with blueprint lem:nonvanishing
\begin{lemma}[nonvanishing of the transforms]\label{lem:nonvanishing}
Assume (A1)--(A3), (A5)--(A7), and let $\Phis{s}{t}f = \mu_{s,t} * f$ be the representation of
Lemma~\ref{lem:convolution-representation}. Then $\hat\mu_{s,t}(\omega) > 0$ for every
$\omega \in \RR$ and every $s \le t$. Consequently
$g_{s,t}(\omega) := -\log\hat\mu_{s,t}(\omega) \in [0,\infty)$ is well defined and even, is
continuous in $\omega$ and continuous in $(s,t)$ on $\{0 \le s \le t\}$, and satisfies
$g_{t,t} = 0$ and $g_{s,t}(0) = 0$.
\end{lemma}

The lemma is elementary but two-axiom: the multiplicativity of (A6) and the
continuity of (A7) are both used, and no clause of Proposition~\ref{prop:fourier-toolbox}
is. Positivity enters only through the hypothesis: the lemma takes the kernels of
Lemma~\ref{lem:convolution-representation} as given, and those are probability measures,
which is what (A4) supplies. Without that the conclusion fails: a signed family of unit
mass can have transforms above $1$ and negative exponents, as the signed compound Poisson
family $e^{h}\sum_n (-h)^n\sigma^{*n}/n!$ with $\sigma = \tfrac12(\delta_{-1}+\delta_1)$ and
$h = t - s$ shows, its transform being $e^{h(1-\cos\omega)}$. What the lemma does not use is
(A4) beyond that. It is the hands-on precursor of the classical fact that an
infinitely divisible law has a zero-free characteristic function (\sscite{sato1999levy},
Lemma~7.5, p.~32), which Theorem~\ref{thm:increments-levy} establishes for the stages.
The order of dependence cannot be reversed, since that theorem's proof
manipulates the exponents this lemma defines.

\begin{proof}
By (A3) and the uniqueness clause of Lemma~\ref{lem:convolution-representation},
$R\mu_{s,t} = \mu_{s,t}$, so $\hat\mu_{s,t}$ is real and even.

\emph{Continuity in $(s,t)$.} Let $\rho$ be the standard Gaussian density, with
$\hat\rho(\omega) = e^{-\omega^2/2} > 0$. If $(s_n,t_n) \to (s,t)$ then
$\mu_{s_n,t_n} * \rho \to \mu_{s,t} * \rho$ in $\Lone$ by (A7), and
\[
  \bigl|\hat\mu_{s_n,t_n}(\omega)\hat\rho(\omega) - \hat\mu_{s,t}(\omega)\hat\rho(\omega)\bigr|
  \ \le\ \|\mu_{s_n,t_n} * \rho - \mu_{s,t} * \rho\|_1 ;
\]
divide by $\hat\rho(\omega)$. The test density must have a zero-free transform, which is why an
indicator --- whose transform vanishes on $2\pi\ZZ \setminus \{0\}$ --- would not do.

\emph{Nonvanishing.} Fix $s$ and $\omega$. The function $t \mapsto \hat\mu_{s,t}(\omega)$ is
continuous on $[s,\infty)$, equals $1$ at $t = s$, and is multiplicative over adjacent intervals
by (A6): $\hat\mu_{s,t} = \hat\mu_{t',t}\,\hat\mu_{s,t'}$ for $s \le t' \le t$. Let
$Z := \{t \ge s : \hat\mu_{s,t}(\omega) = 0\}$, a closed set. If $Z \ne \emptyset$, put
$t_* := \min Z$, which exceeds $s$. For $t' \in (s,t_*)$ we have
$0 = \hat\mu_{s,t_*}(\omega) = \hat\mu_{t',t_*}(\omega)\,\hat\mu_{s,t'}(\omega)$ with the second
factor nonzero by minimality, so $\hat\mu_{t',t_*}(\omega) = 0$; letting $t' \uparrow t_*$,
joint continuity gives $\hat\mu_{t_*,t_*}(\omega) = 0$, against
$\hat\mu_{t_*,t_*} = \hat\delta_0 = 1$. So $Z = \emptyset$: the function is continuous and
zero-free on the interval $[s,\infty)$ and equals $1$ at its left end, hence is positive
throughout, by the intermediate value theorem.

The same conclusion follows without an extremal element, and this is the form the formalisation
takes. Fix $s \le t$. The triangle $\{(a,b) : s \le a \le b \le t\}$ is compact and
$(a,b) \mapsto \hat\mu_{a,b}(\omega)$ is continuous on it, hence uniformly continuous; since the
function is $1$ on the diagonal, there is a single $\delta > 0$ with
$\hat\mu_{a,b}(\omega) > \tfrac12$ whenever $s \le a \le b \le t$ and $b - a < \delta$. Put
$r_n := \min(s + n\delta/2,\, t)$. Then $r_0 = s$, $r_n = t$ for $n$ large, consecutive steps are
shorter than $\delta$, and $\hat\mu_{s,r_{n+1}} = \hat\mu_{r_n,r_{n+1}}\,\hat\mu_{s,r_n}$, so
induction on $n$ carries positivity from $s$ to $t$. Both arguments use exactly the joint
continuity of the first step and nothing else.

\emph{The clauses on $g$.} $g \ge 0$ since $|\hat\mu| \le \|\mu\| = 1$; $g_{s,t}(0) = 0$ by
(A5); continuity in $\omega$ by dominated convergence; evenness and continuity in $(s,t)$ from
the corresponding properties of $\hat\mu$; and $g_{t,t} = 0$ from $\mu_{t,t} = \delta_0$.
\end{proof}

% blueprint rem:modulation-transfer, copied; chapter references made section references
\begin{remark}[modulation transfer; no contrast reversal]\label{rem:modulation-transfer}
On the line the probe $e_\omega$ is a wave and $\omega$ a wavenumber, so
$\hat\mu_{s,t}(\omega) = \Ex\cos(\omega X_{s,t})$ is the stage's \emph{modulation
transfer function}, the contrast a sinusoidal grating of frequency $\omega$ retains after the
random displacement $X_{s,t}$, and the exponent is minus its logarithm.
Lemma~\ref{lem:nonvanishing} then says, in the language of imaging: no stage in a continuous
cascade reverses contrast, none extinguishes a frequency as the box blur does at the zeros of
$\sin\omega/\omega$, and none amplifies. The uniform blur is thereby excluded as a stage of any
family of probability kernels satisfying (A6)--(A7), by an argument that does not yet know
the stages are infinitely divisible. The lemma says nothing about the sign of the \emph{kernel}. The
extended semigroup class of \sscite{pauwels1995extended} beyond $\alpha = 2$ has positive
transfer functions and sign-changing kernels. Those members are excluded by positivity
(A4), through Theorem~\ref{thm:main-characterization}, and not by this lemma.
Transfer functions multiply along the cascade, so optical densities add
(\S\ref{sec:cascade}), and the family is self-similar (\S\ref{sec:covariance}). The reading
of the density's two parts, a Gaussian part and a catalogue of jump displacements, is
Remark~\ref{rem:displacement-dictionary}.
\end{remark}

% blueprint rem:wendel, copied; this section's relation-to-the-causal-theory remark
\begin{remark}[what the representation is, and the relation to the causal theory]\label{rem:wendel}
Wendel's theorem states that the translation-invariant bounded operators on the $\Lone$ of a
locally compact group are exactly the convolutions by finite measures
\sscite{wendel1952left}. Lemma~\ref{lem:convolution-representation} is that step specialised
to $\RR$ and sharpened by (A3)--(A5) to a symmetric probability measure. It is proved rather
than cited. The only imported fact is Proposition~\ref{prop:fourier-uniqueness}. Against the
causal version in \sscite{fagerstrom2026hemigroup}, two devices have fallen away and one has
been added. The clamped exponential is gone, since characters are bounded on all of $\RR$,
which is why \ref{eq:pairing} needs no support hypothesis. The causal support clause is
gone, and with it the one-sided tail bound of the tightness step. Its replacement, that
convolution with a density carried by $[-1,1]$ moves mass by at most $1$, is simpler and
would have served there too. What has been added is Lemma~\ref{lem:nonvanishing}, which on
the half-line was a one-line consequence of the positivity of the integrand. Its continuity
step tests against a Gaussian where the causal proof could test against an indicator.
\end{remark}
